\section{Parse}

\paragraph{Input and output.}
The parse stage consumes Axon source text and produces a one-file Axon AST.
It records syntactic structure, source-level names, declarations, imports,
exports, pragmas, type ascriptions, path literals, and expression forms.

\paragraph{Contract.}
Parsing is intentionally syntax-local.  It does not load imported files,
resolve names, infer types, normalize call syntax, materialize configuration
values, or lower operations.  A parse result may contain unresolved references
and surface sugar; those are obligations of later stages.

\paragraph{Rationale.}
Keeping parse pure makes it possible to test grammar stability independently
from the file system and from checkpoint availability.  It also prevents the
parser from becoming an implicit compatibility layer for downstream stages.

\paragraph{Formal view.}
Let \(\Sigma\) be source text and \(A_f\) be a one-file Axon AST.  Parse is a
partial function
\[
  \mathsf{parse}: \Sigma \rightharpoonup A_f .
\]
It is undefined for syntactically invalid programs, but it is defined for
programs with unresolved identifiers, unresolved imports, or type errors.

\paragraph{Example.}
The source fragment
\begin{verbatim}
import NN

embed :: Tensor[B,S] -> Tensor[B,S,D]
embed ids = NN.embedding@wte ids dim=D
\end{verbatim}
is parsed into an import declaration, a typed definition, a parameter list, and
a call expression whose callee still carries surface path sugar
\code{@wte}.  Parse records this structure but does not decide whether
\code{NN.embedding} exists or whether \code{D} is a valid dimension symbol.
